\documentclass{article}

\usepackage{amsmath,amssymb}
\usepackage{xurl}
\usepackage[hidelinks]{hyperref}
\setlength{\emergencystretch}{2em}

% Shared commands for notes in docs/paper-gaps/.

\DeclareUnicodeCharacter{2080}{\ifmmode{}_{0}\else\textsubscript{0}\fi}
\DeclareUnicodeCharacter{2081}{\ifmmode{}_{1}\else\textsubscript{1}\fi}
\DeclareUnicodeCharacter{2082}{\ifmmode{}_{2}\else\textsubscript{2}\fi}
\DeclareUnicodeCharacter{2099}{\ifmmode{}_{n}\else\textsubscript{n}\fi}

\newcommand{\Lean}{\textsc{Lean}}
\ExplSyntaxOn
\NewDocumentCommand{\leanid}{m}{
  \ifmmode
    \textsf{\detokenize{#1}}
  \else
    \group_begin:
    \urlstyle{sf}
    \tl_set:Nn \l_tmpa_tl {#1}
    \tl_replace_all:Nnn \l_tmpa_tl {\_} {_}
    \exp_args:NV \path \l_tmpa_tl
    \group_end:
  \fi
}
\ExplSyntaxOff
\providecommand{\tr}{\operatorname{tr}}
\newcommand{\tnleanIssueBase}{https://github.com/LionSR/TNLean/issues/}
\newcommand{\tnleanPrBase}{https://github.com/LionSR/TNLean/pull/}
\newcommand{\tnleanDocBase}{https://sirui-lu.com/TNLean/docs/}
\newcommand{\ghissue}[1]{\href{\tnleanIssueBase#1}{GitHub issue \##1}}
\newcommand{\ghpr}[1]{\href{\tnleanPrBase#1}{PR \##1}}
\newcommand{\doclink}[1]{\href{\tnleanDocBase#1.html}{\path{#1}}}

% Dirac notation and the identity operator, as in blueprint/src/macros/common.tex.
% \providecommand keeps note-local definitions authoritative.
\usepackage{dsfont}
\providecommand{\ket}[1]{|#1\rangle}
\providecommand{\bra}[1]{\langle #1|}
\providecommand{\braket}[2]{\langle #1 | #2 \rangle}
\providecommand{\ketbra}[2]{|#1\rangle\!\langle #2|}
\providecommand{\Id}{\mathds{1}}
\providecommand{\one}{\mathds{1}}
\providecommand{\C}{\mathbb{C}}
\providecommand{\MN}[1]{M_{#1}(\C)}
\providecommand{\MD}{M_D(\C)}

% External referencing for paper sources.  The build helper writes sanitized
% auxiliary files under build/paper-aux/ and excludes duplicated source labels.
\usepackage{xr}

% Import a generated paper auxiliary file with a caller-supplied label prefix.
% The candidate paths support builds from docs/paper-gaps/, the repository root,
% and build/paper-aux/ itself.
\newcommand{\importpaperaux}[2]{%
  \IfFileExists{../../build/paper-aux/#2.aux}{%
    \externaldocument[#1]{../../build/paper-aux/#2}%
  }{%
    \IfFileExists{build/paper-aux/#2.aux}{%
      \externaldocument[#1]{build/paper-aux/#2}%
    }{%
      \IfFileExists{#2.aux}{%
        \externaldocument[#1]{#2}%
      }{}%
    }%
  }%
}

% General external reference macros
\newcommand{\paperref}[2]{\ref{#1-#2}}
\newcommand{\papereqref}[2]{(\ref{#1-#2})}

% CPSV16 source (1606.00608 / MPDO-22-12-17-2.tex) external document and shortcuts
\importpaperaux{cpsv16-}{1606.00608/MPDO-22-12-17-2-tnlean}
\newcommand{\cpsvref}[1]{\paperref{cpsv16}{#1}}
\newcommand{\cpsveqref}[1]{\papereqref{cpsv16}{#1}}

% Verdict marker: \gapnote{<kind>}{<status>}. Kind is the policy
% classification (clarification, local-correction, scope-restriction,
% unfaithful, false-source, open-gap); status is open, wip (elimination
% underway), resolved, or historical. Machine-read by the site index;
% prints a verdict line.
\newcommand{\gapnote}[2]{\par\noindent\textbf{Verdict:} #1 (#2).\par}


\title{A Mixed-Prefix Repair of the Two-Site Sector Comparison}
\author{}
\date{2026-08-02}

\begin{document}
\maketitle
\gapnote{local-correction}{resolved}

\section{Source statement and classification}

Appendix~C.4 of Ref.~\cite{Cirac2016MPDO_arXiv} compares the normal tensors in
the one-site and two-site vertical canonical forms. After relabelling, the
comparison at lines~2053--2056 has the form
\begin{align}
    B_\gamma^v
    &=e^{i\theta_\gamma}Z_\gamma A_\gamma^vZ_\gamma^{-1}.
    \label{eq:cpsv16_two_site_sector_unitary_gauge_gap_gauge_phase}
\end{align}
Here $A_\gamma$ is a one-site normal representative, $B_\gamma$ is its
matched two-site representative, and $Z_\gamma$ is invertible. Line~2057 of
Ref.~\cite{Cirac2016MPDO_arXiv} invokes Proposition~4.13 of the same source to
conclude
\begin{align}
    e^{i\theta_\gamma}
    &=1,
    \label{eq:cpsv16_two_site_sector_unitary_gauge_gap_source_phase_removal}\\
    Z_\gamma^\dagger Z_\gamma
    &=\omega_\gamma\mathbf 1,
    \label{eq:cpsv16_two_site_sector_unitary_gauge_gap_source_metric_upgrade}\\
    \omega_\gamma
    &>0.
    \label{eq:cpsv16_two_site_sector_unitary_gauge_gap_source_metric_positivity}
\end{align}
Consequently, $U_\gamma=\omega_\gamma^{-1/2}Z_\gamma$ is unitary and induces
the same conjugation as $Z_\gamma$.

\textbf{Local fix: mixed-prefix alignment.}
The formal proof requires one geometric clarification not stated explicitly
in Ref.~\cite{Cirac2016MPDO_arXiv}. The two-site sector corner must compress
the first two physical sites jointly and leave the remaining sites as
unblocked copies of the original tensor. An ordinary reflected-chain argument
applied to the blocked tensor instead produces a tail of two-site blocks,
which cannot be compared directly with the one-site sector corner. The
mixed-prefix reflection lemma supplies the required common unblocked tail.
This repair changes neither the hypotheses nor the conclusion of
Theorem~4.14(ii)$\Rightarrow$(i) in Ref.~\cite{Cirac2016MPDO_arXiv}; it aligns
the proof at the invocation of Proposition~4.13.

Positivity of two consecutive closed-chain coefficients already removes the
phase in
Eq.~\ref{eq:cpsv16_two_site_sector_unitary_gauge_gap_gauge_phase}. The
remaining issue is the metric identity in
Eq.~\ref{eq:cpsv16_two_site_sector_unitary_gauge_gap_source_metric_upgrade}.

\section{The two corners}

Let $M$ be an MPO tensor that generates matrix product density operators.
Assume that its doubled-index tensor is in literal CPSV canonical form, and
fix a tensor-attached BNT algebra clause. For one matched label, write $A$ for
the one-site normal representative, $B$ for the corresponding two-site
representative, and $Z$ for the exact invertible gauge. Then
\begin{align}
    B^v
    &=ZA^vZ^{-1}.
    \label{eq:cpsv16_two_site_sector_unitary_gauge_gap_exact_gauge}
\end{align}

The one-site vertical coisometric decomposition contains a distinguished
positive copy of $A$. Its inclusion gives a matrix $W$ and a scalar $a>0$
such that
\begin{align}
    W^\dagger W
    &=\mathbf 1,
    \label{eq:cpsv16_two_site_sector_unitary_gauge_gap_one_site_inclusion_isometry}\\
    W^\dagger\widetilde M^vW
    &=aA^v.
    \label{eq:cpsv16_two_site_sector_unitary_gauge_gap_one_site_raw_corner}
\end{align}
This is a genuine one-site physical corner. It uses neither the two-site
gauge nor a Gram identity, fusion datum, or renormalization map.

The matched positive copy in the two-site vertical decomposition gives a
matrix $V$ and a scalar $b>0$ satisfying
\begin{align}
    V^\dagger V
    &=\mathbf 1,
    \label{eq:cpsv16_two_site_sector_unitary_gauge_gap_two_site_inclusion_isometry}\\
    V^\dagger\widetilde M_{[2]}^vV
    &=bZA^vZ^{-1}.
    \label{eq:cpsv16_two_site_sector_unitary_gauge_gap_two_site_gauge_corner}
\end{align}
The formal declarations for these corners are
\begin{center}
\scriptsize
\mbox{\leanid{MPOTensor.BNTAlgebraTensorClause.TwoSiteExactSectorGauge.}}\\
\mbox{\leanid{exists_oneSite_raw_sector_corner},}\qquad
\mbox{\leanid{exists_blockTwo_gauge_sector_corner}.}
\end{center}

\section{The common unblocked reflected tail}

Let $\mathfrak m_N(C;M)$ denote a marked-chain coefficient formed from a
marked tensor $C$ followed by $N$ copies of $M$. The notation suppresses the
open bond indices and the two tail words. Hermiticity of the matrix product
density operators and
Eq.~\ref{eq:cpsv16_two_site_sector_unitary_gauge_gap_one_site_raw_corner}
give
\begin{align}
    \mathfrak m_N(A;M)
    &=\mathfrak m_N(A^{\mathrm{ref}\dagger};M^\dagger).
    \label{eq:cpsv16_two_site_sector_unitary_gauge_gap_raw_reflection}
\end{align}
On the right, Proposition~4.13 of Ref.~\cite{Cirac2016MPDO_arXiv} interchanges
the bond indices and reverses both tail words.

For the two-site corner, compress the first two physical sites jointly
through $V$ and leave the next $N$ sites unblocked. Hermiticity of the same
length-$(N+2)$ density operator gives
\begin{align}
    \mathfrak m_N(\operatorname{gramDressing}(Z,A);M)
    &=\mathfrak m_N(A^{\mathrm{ref}\dagger};M^\dagger).
    \label{eq:cpsv16_two_site_sector_unitary_gauge_gap_dressed_reflection}
\end{align}
The right-hand side is identical to that in
Eq.~\ref{eq:cpsv16_two_site_sector_unitary_gauge_gap_raw_reflection}; in
particular, it has the same unblocked $M^\dagger$ tail. The positive corner
coefficients $a$ and $b$ cancel separately in the two reflection arguments.
Thus, for every positive tail length and all open indices and tail words,
\begin{align}
    \mathfrak m_N(\operatorname{gramDressing}(Z,A);M)
    &=\mathfrak m_N(A;M).
    \label{eq:cpsv16_two_site_sector_unitary_gauge_gap_raw_dressed_marked_chain}
\end{align}

The mixed-prefix identity is formalized by
\begin{center}
\scriptsize
\mbox{\leanid{MPOTensor.IsMPDO.}}\\
\mbox{\leanid{markedChainCoefficient_blockTwoPrefix_}}\\
\mbox{\leanid{gramDressing_eq_reflectedAdjoint}.}
\end{center}
Its essential feature is that the two-site prefix is followed by the original
one-site tensor, not by the blocked tensor.

\section{Both marks lie in the one-site physical-letter span}

Define the positive-definite Gram matrix
\begin{align}
    G
    &=Z^\dagger Z.
    \label{eq:cpsv16_two_site_sector_unitary_gauge_gap_gram_matrix}
\end{align}
The corner in
Eq.~\ref{eq:cpsv16_two_site_sector_unitary_gauge_gap_one_site_raw_corner}
realizes the raw marked tensor in the one-site physical-letter span through
the pair $(W,W)$. Indeed, taking a matrix entry of
$a^{-1}W^\dagger\widetilde M^vW=A^v$ expresses each horizontal slice of $A$
as a linear combination of the physical letters of $M$.

The pair
\begin{align}
    (L_G,R_G)
    &=(WG^\dagger,WG^{-1})
    \label{eq:cpsv16_two_site_sector_unitary_gauge_gap_gram_pair}
\end{align}
realizes the Gram-dressed mark in the same one-site physical-letter span,
because
\begin{align}
    (WG^\dagger)^\dagger\widetilde M^v(WG^{-1})
    &=G(W^\dagger\widetilde M^vW)G^{-1},
    \label{eq:cpsv16_two_site_sector_unitary_gauge_gap_gram_compression_first}\\
    &=aGA^vG^{-1},
    \label{eq:cpsv16_two_site_sector_unitary_gauge_gap_gram_compression_second}\\
    &=a\operatorname{gramDressing}(Z,A)^v.
    \label{eq:cpsv16_two_site_sector_unitary_gauge_gap_gram_one_site_compression}
\end{align}
Both sides of
Eq.~\ref{eq:cpsv16_two_site_sector_unitary_gauge_gap_raw_dressed_marked_chain}
therefore satisfy the physical-letter premise of the literal CPSV Lemma~L in
Ref.~\cite{Cirac2016MPDO_arXiv}. This is the second essential part of the
repair: no identity-gauge corner is needed inside the blocked physical space.

Apply Lemma~L of Ref.~\cite{Cirac2016MPDO_arXiv} to
Eq.~\ref{eq:cpsv16_two_site_sector_unitary_gauge_gap_raw_dressed_marked_chain}.
Literal CPSV canonical form supplies representative separation at a
sufficiently large positive tail length, so
\begin{align}
    \operatorname{gramDressing}(Z,A)
    &=A.
    \label{eq:cpsv16_two_site_sector_unitary_gauge_gap_gram_dressing_target}
\end{align}
Equivalently,
\begin{align}
    GA^vG^{-1}
    &=A^v
    \quad\text{for every vertical letter }v.
    \label{eq:cpsv16_two_site_sector_unitary_gauge_gap_gram_commutation}
\end{align}

Since $A$ is normal, its letters generate the full matrix algebra after a
finite blocking. Equation~\ref{eq:cpsv16_two_site_sector_unitary_gauge_gap_gram_commutation}
places $G$ in the scalar commutant. Since $G=Z^\dagger Z$ is positive
definite, a real number $\omega>0$ exists such that
\begin{align}
    Z^\dagger Z
    &=\omega\mathbf 1.
    \label{eq:cpsv16_two_site_sector_unitary_gauge_gap_gram_scalar}
\end{align}
Define $U=\omega^{-1/2}Z$. Then
\begin{align}
    U
    &=\omega^{-1/2}Z,
    \label{eq:cpsv16_two_site_sector_unitary_gauge_gap_unitary_definition}\\
    U^\dagger U
    &=\mathbf 1,
    \label{eq:cpsv16_two_site_sector_unitary_gauge_gap_unitary_normalization}\\
    B^v
    &=UA^vU^\dagger.
    \label{eq:cpsv16_two_site_sector_unitary_gauge_gap_unitary_gauge}
\end{align}
This proves the unitary normalization at line~2057 of
Ref.~\cite{Cirac2016MPDO_arXiv} without assuming a positive blocked raw
corner, fusion coisometries, or renormalization maps.

\section{Completion of Theorem~4.14}

Applying the construction sector by sector gives simultaneous unitaries for
the one-site and two-site vertical decompositions. The direct-sum
conjugations, vertical retractions, and normalized restorations then produce
the two physical maps written in Appendix~C.4, lines~2065--2085, of
Ref.~\cite{Cirac2016MPDO_arXiv}. The trace-restoring completion on discarded
zero-sector complements is the local fix documented in
\path{docs/paper-gaps/cpgsv17_vertical_isometry_zero_sector.tex}; its added
terms vanish on the physical closures generated by $M$.

A tensor-attached BNT algebra clause, together with the standing literal CPSV
canonical-form and MPDO hypotheses, therefore implies the trace-preserving
completely positive renormalization identities of Definition~4.1 in
Ref.~\cite{Cirac2016MPDO_arXiv}. This is the unconditional implication
Theorem~4.14(ii)$\Rightarrow$(i), formalized by
\begin{center}
\scriptsize
\leanid{MPOTensor.BNTAlgebraTensorClause.isRFPViaTS}.
\end{center}

The formalized forward implication produces the corrected active-support
fusion clause, and active-support fusion implies the tensor-attached algebra
clause. Combining these directions gives
\begin{align}
    \operatorname{IsRFPViaTS}(M)
    &\Longleftrightarrow
      \operatorname{HasBNTAlgebraTensorClause}(M),
    \label{eq:cpsv16_two_site_sector_unitary_gauge_gap_rfp_algebra_equivalence}\\
    \operatorname{IsRFPViaTS}(M)
    &\Longleftrightarrow
      \operatorname{HasBNTFusionTensorClause}(M).
    \label{eq:cpsv16_two_site_sector_unitary_gauge_gap_rfp_fusion_equivalence}
\end{align}
These equivalences hold under exactly the standing canonical-form and MPDO
hypotheses of Theorem~4.14 in Ref.~\cite{Cirac2016MPDO_arXiv}. The second
equivalence uses the corrected active-support interpretation of
statement~(iii), including the documented empty-support and coisometry local
fixes. Neither equivalence assumes a stronger horizontal canonical form or an
additional reflected-target hypothesis. The corresponding formal
declarations are
\begin{center}
\scriptsize
\mbox{\leanid{MPOTensor.isRFPViaTS_iff_hasBNTAlgebraTensorClause},}\\
\mbox{\leanid{MPOTensor.isRFPViaTS_iff_hasBNTFusionTensorClause}.}
\end{center}

\section{Why the former direct route stopped}

The mixed-prefix proof does not invalidate the earlier obstruction analyses.
Those analyses concern a stronger and unnecessary route: constructing a raw
representative as a same-sided corner of the blocked physical tensor and then
comparing two ordinary blocked chains.

\subsection{The circular blocked raw corner}

From
Eq.~\ref{eq:cpsv16_two_site_sector_unitary_gauge_gap_two_site_gauge_corner},
the natural candidate for a blocked physical inclusion of $A$ is a scalar
multiple of $VZ$. It is an isometry exactly when
$Z^\dagger Z=\omega\mathbf 1$. Once this identity is known,
$\omega^{-1/2}VZ$ gives a positive same-sided raw corner. Conversely, a raw
isometric corner of this form, compared with $V$, proves the same Gram
identity. The construction is therefore circular if used to establish
Eq.~\ref{eq:cpsv16_two_site_sector_unitary_gauge_gap_gram_scalar}.

The repair avoids this circle. It uses the independent one-site corner $W$
for the raw marked chain and uses $V$ only for the two-site prefix. The two
corners need not lie in the same physical space because their reflection
identities have the same unblocked tail.

\subsection{The oblique blocked compression}

The invertible gauge always gives the oblique maps
\begin{align}
    L
    &=V(Z^{-1})^\dagger,
    \label{eq:cpsv16_two_site_sector_unitary_gauge_gap_oblique_left_map}\\
    R
    &=VZ.
    \label{eq:cpsv16_two_site_sector_unitary_gauge_gap_oblique_right_map}
\end{align}
They satisfy
\begin{align}
    L^\dagger\widetilde M_{[2]}^vR
    &=bA^v,
    \label{eq:cpsv16_two_site_sector_unitary_gauge_gap_oblique_raw_compression}\\
    R^\dagger\widetilde M_{[2]}^vL
    &=b(Z^\dagger Z)A^v(Z^\dagger Z)^{-1}.
    \label{eq:cpsv16_two_site_sector_unitary_gauge_gap_oblique_gram_compression}
\end{align}
This is a useful physical-letter realization, but it is not a positive
same-sided corner. Adjoint reflection exchanges $L$ and $R$, so the opposite
orientation is Gram-dressed. The oblique blocked construction alone therefore
does not prove
Eq.~\ref{eq:cpsv16_two_site_sector_unitary_gauge_gap_gram_dressing_target}.
The mixed-prefix proof neither makes this false inference nor converts $L$
and $R$ into one map.

\subsection{Fusion data cannot be imported into the converse}

A local coisometry $F_{\alpha\beta}$ satisfying
\begin{align}
    F_{\alpha\beta}(A_\alpha A_\beta)^vF_{\alpha\beta}^\dagger
    &=\bigoplus_\gamma
      \chi_{\alpha\beta\gamma}\otimes A_\gamma^v
    \label{eq:cpsv16_two_site_sector_unitary_gauge_gap_fusion_compression}
\end{align}
would immediately produce isometric blocked raw corners. Such an
$F_{\alpha\beta}$ is precisely the fusion datum in statement~(iii) of
Theorem~4.14 in Ref.~\cite{Cirac2016MPDO_arXiv}. Using it to prove
statement~(ii)$\Rightarrow$(i) would be circular. The algebra clause supplies
the closed-chain product law, not the local map in
Eq.~\ref{eq:cpsv16_two_site_sector_unitary_gauge_gap_fusion_compression}. The
completed converse proof does not use fusion data.

\section{Counterexamples delimiting the direct route}

\subsection{Closed chains and an invertible gauge are insufficient}

Let the letters of $A$ include the identity and the matrix units of
$M_2(\mathbb C)$. Define
\begin{align}
    X
    &=\begin{pmatrix}2&0\\0&1\end{pmatrix},
    \label{eq:cpsv16_two_site_sector_unitary_gauge_gap_generic_similarity_matrix}\\
    C^v
    &=XA^vX^{-1}.
    \label{eq:cpsv16_two_site_sector_unitary_gauge_gap_generic_similarity_tensor}
\end{align}
The tensors $A$ and $C$ have identical closed chains at every positive
length. Suppose this fact alone produced a square isometry $Q$ and a positive
scalar $c$ such that $Q^\dagger C^vQ=cA^v$. The identity letter would force
$c=1$. Then $Q$ would be unitary, and $Q^\dagger X$ would commute with every
matrix and hence be scalar. This would force $X^\dagger X$ to be scalar,
contrary to
\begin{align}
    X^\dagger X
    &=\begin{pmatrix}4&0\\0&1\end{pmatrix}.
    \label{eq:cpsv16_two_site_sector_unitary_gauge_gap_generic_similarity_gram}
\end{align}
Likewise, for a letter $A^{v_0}=E_{01}$,
\begin{align}
    \operatorname{gramDressing}(X,A)^{v_0}
    &=4E_{01},
    \label{eq:cpsv16_two_site_sector_unitary_gauge_gap_generic_gram_dressed_letter}\\
    4E_{01}
    &\neq E_{01}=A^{v_0}.
    \label{eq:cpsv16_two_site_sector_unitary_gauge_gap_generic_gram_obstruction}
\end{align}

These examples obstruct any derivation of a blocked raw corner or Gram
identity from positive-length closed-chain equality and an invertible gauge
alone. They do not contradict the completed tensor-attached theorem because
they do not supply one MPO tensor satisfying the algebra clause, literal CPSV
canonical form, and MPDO positivity from which both physical corners arise.
The relevant formal declarations are
\begin{center}
\scriptsize
\mbox{\leanid{MPSTensor.PositiveMinimalRealizationCounterexample.}}\\
\mbox{\leanid{sameMPV}\ensuremath{{}_{2}}\leanid{Pos_does_not_force_positive_}}\\
\mbox{\leanid{isometric_realization},}\qquad
\mbox{\leanid{gramDressing_gauge_ne_one}.}
\end{center}

The same distinction explains the negative result for the per-block linear
extension. A positive algebra automorphism of a full matrix algebra is a
star automorphism and therefore has a unitary implementer by unitary
Skolem--Noether. The generic nonunitary similarity makes the corresponding
per-block extension nonpositive. Positivity of that extension therefore
cannot be inferred from closed chains alone. The mixed-prefix proof does not
require this positivity. The corresponding formal declarations are
\begin{center}
\scriptsize
\mbox{\leanid{Matrix.exists_unitary_conj_of_positive_algEquiv_matrix},}\\
\mbox{\leanid{MPSTensor.exists_unitary_conj_of_positive_perBlockLinearExtension},}\\
\mbox{\leanid{MPSTensor.PositiveMinimalRealizationCounterexample.}}\\
\mbox{\leanid{perBlockLinearExtension_not_positive}.}
\end{center}

\subsection{Terminal positivity and the bond-two test}

For every retained one-site sector, tensor-attached MPDO positivity implies
\begin{align}
    T_\gamma
    &=\sum_a A_\gamma^{(a,a)},
    \label{eq:cpsv16_two_site_sector_unitary_gauge_gap_terminal_sector_transfer_definition}\\
    T_\gamma
    &\succeq0.
    \label{eq:cpsv16_two_site_sector_unitary_gauge_gap_terminal_sector_transfer_positivity}
\end{align}
This condition excludes the simplest nonunitary minimal-realization witness
from the full theorem because its terminal transfer is not Hermitian. The
declaration
\begin{center}
\scriptsize
\leanid{MPOTensor.BNTAlgebraTensorClause.%
physTraceTransfer_verticalBNTMPO_posSemidef}
\end{center}
formalizes this necessary condition.

The bond-two singleton family gives a sharper diagnostic. Its undeformed
member is an all-length MPDO with a genuine one-label tensor-attached algebra
clause. A concrete nonunitary physical similarity preserves the algebra
clause and literal horizontal CPSV canonical form, but its length-two MPO is
not positive. It therefore fails the MPDO hypothesis. This example shows that
the positivity hypothesis is necessary; it is not a counterexample to
Theorem~4.14 of Ref.~\cite{Cirac2016MPDO_arXiv}. Its former role as evidence
for a metric boundary is superseded by the mixed-prefix comparison. The
formal tests are
\begin{center}
\scriptsize
\mbox{\leanid{MPOTensor.BondTwoSingletonBaseModel.}}\\
\mbox{\leanid{gaugeGram_eq_pos_smul_one_of_gaugeDeformedBaseMPO_isMPDO},}\\
\mbox{\leanid{gaugeDeformedBaseMPO_gauge_not_isMPDO},}\\
\mbox{\leanid{gaugeDeformedBaseMPOBNTAlgebraTensorClause},}\qquad
\mbox{\leanid{gaugeDeformedBaseMPO_toMPSTensor_isCPSVCanonicalForm}.}
\end{center}

\section{Verdict}

The two-site sector comparison no longer stops at an invertible gauge. The
one-site raw corner and the two-site gauge corner satisfy reflection
identities with the same unblocked tail. Their common reflected target equates
the raw and Gram-dressed marked chains. The pairs $(W,W)$ and
$(WG^\dagger,WG^{-1})$ place both marks in the one-site physical-letter span.
Literal CPSV Lemma~L in Ref.~\cite{Cirac2016MPDO_arXiv} then gives
$\operatorname{gramDressing}(Z,A)=A$. Normality implies
$Z^\dagger Z=\omega\mathbf 1$ with $\omega>0$, after which unitary
normalization and the physical renormalization maps follow.

This is a local repair of the source proof, not a change to its mathematical
statement. It discharges the comparison formerly tracked in issues 5284,
4648, and 4645, establishes the unconditional conclusion at line~2057 of
Ref.~\cite{Cirac2016MPDO_arXiv}, and completes both the tensor-attached
algebra equivalence and the corrected active-support fusion equivalence
associated with Theorem~4.14. The direct-route counterexamples remain useful:
closed chains and invertible similarity do not by themselves produce a
blocked raw corner. The completed proof succeeds because it compares a
one-site corner with a mixed two-site prefix inside the same MPDO family.

\bibliographystyle{alpha}
\bibliography{references}

\end{document}
